% Practice list A --- topics of Listas 1-2 (Kurlat chs. 1-5).
% \solucoesfalse -> blank list ; \solucoestrue -> answer key + rubric.
% Compile: pdflatex (2x); check: pdffonts (all Type 1). Numbers: practice-lists-check.py
\documentclass[11pt,a4paper]{article}
% Shared preamble of the practice lists (practice-L*.tex), adapted from lista-exam.tex.
% Each list sets \solucoesfalse / \solucoestrue and \listatema right after % Shared preamble of the practice lists (practice-L*.tex), adapted from lista-exam.tex.
% Each list sets \solucoesfalse / \solucoestrue and \listatema right after % Shared preamble of the practice lists (practice-L*.tex), adapted from lista-exam.tex.
% Each list sets \solucoesfalse / \solucoestrue and \listatema right after \input{practice-preamble}.
\usepackage[utf8]{inputenc}
\usepackage[T1]{fontenc}
\usepackage{lmodern}   % Type 1 vector fonts (replaces PK bitmaps)
\usepackage{cmap}      % ToUnicode CMap -> copyable text and accents
\usepackage[english]{babel}
\usepackage{amsmath,amssymb,amsfonts}
\usepackage[a4paper,margin=2.2cm]{geometry}
\usepackage{enumitem}
\usepackage{fancyhdr}
\usepackage{comment}
\usepackage{xcolor}
\usepackage{microtype}
\usepackage{booktabs}
\usepackage{tikz}

\newif\ifsolucoes
\newcommand{\disciplina}{Macroeconomics I --- MPE Insper}
\newcommand{\espacoresol}{4.5cm}   % height of the answer space in the blank list
% The sol environment must be decided after the list sets its toggle.
\AtBeginDocument{\ifsolucoes\includecomment{sol}\else\excludecomment{sol}\fi}

\pagestyle{fancy}\fancyhf{}
\lhead{\footnotesize\disciplina}
\rhead{\footnotesize Practice list\ifsolucoes{} --- Solutions and rubric\fi}
\cfoot{\footnotesize\thepage}
\renewcommand{\headrulewidth}{0.3pt}
\setlength{\headheight}{14pt}
\setlist{topsep=2pt,itemsep=2pt}
\newcounter{questao}

\newcommand{\bloco}[3]{\bigskip\noindent
  {\large\bfseries\color{blue!55!black}Block #1 --- #2}\hfill\textit{(#3)}\par
  \rule{\linewidth}{0.3pt}\smallskip}
\newcommand{\questao}[1]{\refstepcounter{questao}\medskip\noindent
  \textbf{Question \thequestao.}~#1\par}
\newenvironment{alternativas}{\begin{enumerate}[label=\Alph*.,leftmargin=2.2em,topsep=1pt]}{\end{enumerate}}
\newcommand{\gabarito}[1]{\ifsolucoes\par\noindent\textbf{Answer:}~#1\fi}
\newcommand{\resolespaco}{\ifsolucoes\else\par\textit{Solution:}\vspace{\espacoresol}\par\fi}
\newcommand{\passo}[2]{\par\smallskip\noindent\textbf{Step --- #1.}~#2}
\newcommand{\intuicao}[1]{\par\smallskip\noindent\textit{Economic intuition.}~#1}
\newcommand{\refbib}[1]{\par\smallskip\noindent{\footnotesize\color{black!60}\textbf{Ref:} #1}}
\newcommand{\rubrica}[1]{\par\smallskip\noindent\textbf{Grading rubric.}\begin{itemize}[leftmargin=1.4em,topsep=1pt]#1\end{itemize}}

\newcommand{\cabecalho}{%
\begin{center}
  {\Large\bfseries \disciplina}\\[2pt]
  {\large Practice list: \listatema}\ifsolucoes\\[2pt]\textit{Solutions and grading rubric}\fi
\end{center}
\vspace{2pt}
{\footnotesize\noindent\textbf{Instructions.} Multiple choice: mark the single correct option (no
penalty for wrong answers). \emph{True/False}: every 2 wrong items cancel 1 right one
(score $=\max\{0,\text{right}-\text{wrong}/2\}\times$ item value); blank items do not count.
Discursive questions: show every step. Numbers are illustrative and chosen to give clean answers.\par}}
.
\usepackage[utf8]{inputenc}
\usepackage[T1]{fontenc}
\usepackage{lmodern}   % Type 1 vector fonts (replaces PK bitmaps)
\usepackage{cmap}      % ToUnicode CMap -> copyable text and accents
\usepackage[english]{babel}
\usepackage{amsmath,amssymb,amsfonts}
\usepackage[a4paper,margin=2.2cm]{geometry}
\usepackage{enumitem}
\usepackage{fancyhdr}
\usepackage{comment}
\usepackage{xcolor}
\usepackage{microtype}
\usepackage{booktabs}
\usepackage{tikz}

\newif\ifsolucoes
\newcommand{\disciplina}{Macroeconomics I --- MPE Insper}
\newcommand{\espacoresol}{4.5cm}   % height of the answer space in the blank list
% The sol environment must be decided after the list sets its toggle.
\AtBeginDocument{\ifsolucoes\includecomment{sol}\else\excludecomment{sol}\fi}

\pagestyle{fancy}\fancyhf{}
\lhead{\footnotesize\disciplina}
\rhead{\footnotesize Practice list\ifsolucoes{} --- Solutions and rubric\fi}
\cfoot{\footnotesize\thepage}
\renewcommand{\headrulewidth}{0.3pt}
\setlength{\headheight}{14pt}
\setlist{topsep=2pt,itemsep=2pt}
\newcounter{questao}

\newcommand{\bloco}[3]{\bigskip\noindent
  {\large\bfseries\color{blue!55!black}Block #1 --- #2}\hfill\textit{(#3)}\par
  \rule{\linewidth}{0.3pt}\smallskip}
\newcommand{\questao}[1]{\refstepcounter{questao}\medskip\noindent
  \textbf{Question \thequestao.}~#1\par}
\newenvironment{alternativas}{\begin{enumerate}[label=\Alph*.,leftmargin=2.2em,topsep=1pt]}{\end{enumerate}}
\newcommand{\gabarito}[1]{\ifsolucoes\par\noindent\textbf{Answer:}~#1\fi}
\newcommand{\resolespaco}{\ifsolucoes\else\par\textit{Solution:}\vspace{\espacoresol}\par\fi}
\newcommand{\passo}[2]{\par\smallskip\noindent\textbf{Step --- #1.}~#2}
\newcommand{\intuicao}[1]{\par\smallskip\noindent\textit{Economic intuition.}~#1}
\newcommand{\refbib}[1]{\par\smallskip\noindent{\footnotesize\color{black!60}\textbf{Ref:} #1}}
\newcommand{\rubrica}[1]{\par\smallskip\noindent\textbf{Grading rubric.}\begin{itemize}[leftmargin=1.4em,topsep=1pt]#1\end{itemize}}

\newcommand{\cabecalho}{%
\begin{center}
  {\Large\bfseries \disciplina}\\[2pt]
  {\large Practice list: \listatema}\ifsolucoes\\[2pt]\textit{Solutions and grading rubric}\fi
\end{center}
\vspace{2pt}
{\footnotesize\noindent\textbf{Instructions.} Multiple choice: mark the single correct option (no
penalty for wrong answers). \emph{True/False}: every 2 wrong items cancel 1 right one
(score $=\max\{0,\text{right}-\text{wrong}/2\}\times$ item value); blank items do not count.
Discursive questions: show every step. Numbers are illustrative and chosen to give clean answers.\par}}
.
\usepackage[utf8]{inputenc}
\usepackage[T1]{fontenc}
\usepackage{lmodern}   % Type 1 vector fonts (replaces PK bitmaps)
\usepackage{cmap}      % ToUnicode CMap -> copyable text and accents
\usepackage[english]{babel}
\usepackage{amsmath,amssymb,amsfonts}
\usepackage[a4paper,margin=2.2cm]{geometry}
\usepackage{enumitem}
\usepackage{fancyhdr}
\usepackage{comment}
\usepackage{xcolor}
\usepackage{microtype}
\usepackage{booktabs}
\usepackage{tikz}

\newif\ifsolucoes
\newcommand{\disciplina}{Macroeconomics I --- MPE Insper}
\newcommand{\espacoresol}{4.5cm}   % height of the answer space in the blank list
% The sol environment must be decided after the list sets its toggle.
\AtBeginDocument{\ifsolucoes\includecomment{sol}\else\excludecomment{sol}\fi}

\pagestyle{fancy}\fancyhf{}
\lhead{\footnotesize\disciplina}
\rhead{\footnotesize Practice list\ifsolucoes{} --- Solutions and rubric\fi}
\cfoot{\footnotesize\thepage}
\renewcommand{\headrulewidth}{0.3pt}
\setlength{\headheight}{14pt}
\setlist{topsep=2pt,itemsep=2pt}
\newcounter{questao}

\newcommand{\bloco}[3]{\bigskip\noindent
  {\large\bfseries\color{blue!55!black}Block #1 --- #2}\hfill\textit{(#3)}\par
  \rule{\linewidth}{0.3pt}\smallskip}
\newcommand{\questao}[1]{\refstepcounter{questao}\medskip\noindent
  \textbf{Question \thequestao.}~#1\par}
\newenvironment{alternativas}{\begin{enumerate}[label=\Alph*.,leftmargin=2.2em,topsep=1pt]}{\end{enumerate}}
\newcommand{\gabarito}[1]{\ifsolucoes\par\noindent\textbf{Answer:}~#1\fi}
\newcommand{\resolespaco}{\ifsolucoes\else\par\textit{Solution:}\vspace{\espacoresol}\par\fi}
\newcommand{\passo}[2]{\par\smallskip\noindent\textbf{Step --- #1.}~#2}
\newcommand{\intuicao}[1]{\par\smallskip\noindent\textit{Economic intuition.}~#1}
\newcommand{\refbib}[1]{\par\smallskip\noindent{\footnotesize\color{black!60}\textbf{Ref:} #1}}
\newcommand{\rubrica}[1]{\par\smallskip\noindent\textbf{Grading rubric.}\begin{itemize}[leftmargin=1.4em,topsep=1pt]#1\end{itemize}}

\newcommand{\cabecalho}{%
\begin{center}
  {\Large\bfseries \disciplina}\\[2pt]
  {\large Practice list: \listatema}\ifsolucoes\\[2pt]\textit{Solutions and grading rubric}\fi
\end{center}
\vspace{2pt}
{\footnotesize\noindent\textbf{Instructions.} Multiple choice: mark the single correct option (no
penalty for wrong answers). \emph{True/False}: every 2 wrong items cancel 1 right one
(score $=\max\{0,\text{right}-\text{wrong}/2\}\times$ item value); blank items do not count.
Discursive questions: show every step. Numbers are illustrative and chosen to give clean answers.\par}}

\solucoesfalse                 % <<<< \solucoestrue for the ANSWER KEY
\newcommand{\listatema}{Measurement, Solow and Growth Accounting (Listas 1--2)}

\begin{document}
\cabecalho

% =============================================================================
\bloco{1}{Measuring output, prices and welfare}{Kurlat chs.~1--2 \,\textbullet\, 25 points}

\questao{(10 points) Objective items. IBGE's national accounts team tests its procedures on a
two-good economy (food, good 1; electronics, good 2):}

\begin{center}
\begin{tabular}{lcccc}
\toprule
 & $p_1$ & $q_1$ & $p_2$ & $q_2$ \\
\midrule
2024 & 2 & 10 & 10 & 5 \\
2025 & 3 & 12 & 7 & 8 \\
\bottomrule
\end{tabular}
\end{center}

\textbf{a)} (3 points) With 2024 as the base year, the GDP deflator for 2025 (2024 $=100$) is:
\begin{alternativas}
\item $92.86$
\item $131.43$
\item $88.46$
\item $113.04$
\end{alternativas}
\gabarito{C}
\begin{sol}
\passo{nominal and real}{Nominal GDP 2025 $=3\cdot12+7\cdot8=92$. Real GDP 2025 at 2024 prices
$=2\cdot12+10\cdot8=104$. Deflator $=100\cdot 92/104=88.46$.}
\passo{distractors}{$92.86=65/70$ is the Laspeyres (fixed 2024 basket) price index, not the
deflator; $131.43=92/70$ is nominal GDP growth; $113.04$ inverts the ratio.}
\refbib{Kurlat (2020), \S1.2, p.~22.}
\end{sol}

\medskip
\textbf{b)} (3 points) Real GDP growth between 2024 and 2025 computed with Kurlat's chained
method (geometric average of the growth rates at initial-year and final-year prices) is:
\begin{alternativas}
\item $48.57\%$
\item $31.43\%$
\item $41.54\%$
\item $45.01\%$
\end{alternativas}
\gabarito{D}
\begin{sol}
\passo{two growth rates}{At 2024 prices: $1+g^I=104/70=1.4857$. At 2025 prices: real 2024
$=3\cdot10+7\cdot5=65$, so $1+g^F=92/65=1.4154$.}
\passo{chain}{$1+g=(1.4857\cdot1.4154)^{1/2}=1.4501$, i.e.\ $45.01\%$. Option B is nominal growth.}
\intuicao{Electronics grows most and its relative price falls, so early-year prices weight the
fast-growing good more and give the higher rate; chaining splits the difference.}
\refbib{Kurlat (2020), \S1.2, ``Alternative 3: chained prices'', pp.~23--24.}
\end{sol}

\medskip
\textbf{c)} (2 points, T or F) If the IPCA, a Laspeyres index with a fixed basket, rises 6\% in a
year, then the extra income a household needs to reach its original utility level rose by
\emph{at least} 6\%.
\gabarito{F}
\begin{sol}
\passo{substitution bias}{The fixed basket still affords the old utility, so 6\% is enough; the
household can do better by substituting toward goods that became relatively cheaper. The
Laspeyres index is an \emph{upper} bound on the true cost-of-living change: at most 6\%.}
\refbib{Kurlat (2020), \S1.2; formula sheet \texttt{rules/01\_mensuracao\_agregados.md}.}
\end{sol}

\medskip
\textbf{d)} (2 points, T or F) Suppose tradable goods obey the law of one price at the market
exchange rate and non-tradable services are cheaper in Brazil than in the US. Then converting
Brazil's GDP to dollars at the market exchange rate understates it relative to a PPP conversion.
\gabarito{T}
\begin{sol}
\passo{why}{At the market rate tradables are priced the same in both countries, but Brazil's
non-tradables are valued at their (low) Brazilian price. PPP values them at common prices, so
Brazilian output is worth more at PPP (Balassa--Samuelson).}
\refbib{Kurlat (2020), \S1.2, ``Comparisons across countries and PPP''.}
\end{sol}

% -----------------------------------------------------------------------------
\questao{(15 points) Discursive. A Penn World Table analyst compares Brazil and the US. Each
produces a tradable good $T$ and a non-tradable service $N$. Prices are in local currency;
the market exchange rate is $e=5$ reais per dollar.}

\begin{center}
\begin{tabular}{lcccc}
\toprule
 & $p_T$ & $q_T$ & $p_N$ & $q_N$ \\
\midrule
Brazil (reais) & 10 & 50 & 20 & 100 \\
US (dollars) & 2 & 200 & 20 & 150 \\
\bottomrule
\end{tabular}
\end{center}

\textbf{a)} (5 points) Compute Brazil's GDP in dollars (i) at the market exchange rate and
(ii) at PPP, using US prices as the common prices. Compute Brazil/US GDP under each
measure and the implied PPP exchange rate. Explain the gap.
\resolespaco
\begin{sol}
\passo{local GDP}{Brazil: $10\cdot50+20\cdot100=2500$ reais. US: $2\cdot200+20\cdot150=3400$ dollars.}
\passo{market rate}{$2500/5=500$ dollars; Brazil/US $=500/3400=0.147$.}
\passo{PPP}{Brazil's quantities at US prices: $2\cdot50+20\cdot100=2100$ dollars;
Brazil/US $=2100/3400=0.618$. Implied PPP rate: $2500/2100=1.19$ reais per dollar.}
\passo{gap}{Tradables obey the law of one price ($10/2=5=e$), but the service costs
$20/5=4$ dollars in Brazil against 20 in the US. The market rate prices Brazil's large service
sector at a fifth of its PPP value, so Brazil looks 4.2 times poorer than it is.}
\intuicao{Exchange rates are set by traded goods; living standards depend on all goods. Poorer
countries have cheap non-tradables (low wages in services), so the market rate understates them.}
\refbib{Kurlat (2020), \S1.2, pp.~24--27.}
\rubrica{
\item Local GDP of both countries: 1 pt.
\item Market-rate conversion and ratio: 1 pt.
\item PPP value, ratio and implied PPP rate: 2 pts.
\item Explanation through non-tradables / law of one price: 1 pt. Saying only ``PPP is better''
  with no mechanism: 0 for this line.
\item Converting PPP by dividing by $e$ again (double conversion): max.\ 2 pts in the item.
\item Arithmetic slip with correct method: $-1$ pt.
}
\end{sol}

\textbf{b)} (5 points) Now use the Jones--Klenow welfare measure of Kurlat \S2.2 with $\sigma=1$
and equal leisure in both countries, so that a random person's utility is
$\mathbb{E}\big[a(\bar u+\ln c)\big]$, where $a=1$ if alive. Normalize $c_{US}=1$, set
$\bar u=5$, survival probabilities $a_{US}=0.79$, $a_{BR}=0.76$, and Brazilian consumption per
person $\bar c_{BR}=0.30$, \emph{equal for everyone}. Derive $\lambda$ (the scaling of US
consumption that leaves Rawls indifferent), compute it, and split the welfare gap into a
consumption part and a life-expectancy part.
\resolespaco
\begin{sol}
\passo{indifference}{$a_{US}\,[\bar u+\ln(\lambda c_{US})]=a_{BR}\,[\bar u+\ln \bar c_{BR}]$, so
$\ln\lambda=\dfrac{a_{BR}}{a_{US}}(\bar u+\ln \bar c_{BR})-\bar u$.}
\passo{number}{$\ln\lambda=\frac{0.76}{0.79}(5-1.204)-5=-1.348$, so $\lambda=0.260$.}
\passo{decomposition}{With equal survival $\lambda$ would be $0.30$. Lower life expectancy
multiplies it by $0.260/0.30=0.866$: about 13\% of consumption-equivalent welfare.}
\intuicao{$\bar u$ is the value of being alive. Three fewer years of life expectancy cost Rawls
the chance of enjoying $\bar u+\ln c$, so he needs less US consumption to be indifferent.}
\refbib{Kurlat (2020), \S2.2, eqs.\ (2.2.1)--(2.2.2), pp.~33--41.}
\rubrica{
\item Writes the indifference condition with $a$ multiplying the whole bracket: 2 pts.
\item Solves for $\ln\lambda$ and computes $\lambda\approx0.26$: 2 pts.
\item Decomposition and interpretation: 1 pt.
\item Multiplying only $\ln c$ by $a$ (forgetting $\bar u$): max.\ 2 pts in the item.
}
\end{sol}

\textbf{c)} (5 points) Now half of Brazilians consume $0.12$ and half $0.48$ (same mean $0.30$).
Recompute $\lambda$. By how much does inequality lower it? Would a higher $\sigma$ make the
penalty larger or smaller? Why does Brazil fall further behind in welfare than in consumption?
\resolespaco
\begin{sol}
\passo{expected log}{$\mathbb{E}\ln c=\tfrac12\ln0.12+\tfrac12\ln0.48=\ln\sqrt{0.0576}=\ln 0.24$.}
\passo{new $\lambda$}{$\ln\lambda=\frac{0.76}{0.79}(5+\ln0.24)-5=-1.563$, so $\lambda=0.210$.}
\passo{penalty}{$0.210/0.260=0.807$: inequality removes about 19\%. It acts through the ratio
geometric/arithmetic mean $0.24/0.30=0.8$, scaled by survival ($0.8^{0.96}$).}
\passo{role of $\sigma$}{$\sigma$ is the curvature of $u$. With $\sigma>1$, utility is more concave
and the poor half's low consumption weighs more: the penalty grows.}
\intuicao{Behind the veil of ignorance, inequality is risk. A risk-averse Rawls pays to avoid it,
so an unequal country with the same mean consumption is worth less to him.}
\refbib{Kurlat (2020), \S2.2, risk aversion and eq.\ (2.2.5), pp.~35--38.}
\rubrica{
\item Uses $\mathbb{E}\ln c$, not $\ln\mathbb{E}c$: 2 pts. Using $\ln 0.30$: max.\ 1 pt in the item.
\item New $\lambda\approx0.21$ and the penalty: 2 pts.
\item Direction of $\sigma$ with the concavity argument: 1 pt.
}
\end{sol}

% =============================================================================
\bloco{2}{Solow, the Golden Rule and growth accounting}{Kurlat chs.~4--5 \,\textbullet\, 25 points}

\questao{(10 points) Objective items. IPEA calibrates a Solow model \emph{per worker}, without
technological progress: $y=k^{\alpha}$, $\alpha=1/3$, $s=0.20$, $n=0.02$, $\delta=0.08$.}

\textbf{a)} (3 points) Steady-state output per worker $y^*$ is:
\begin{alternativas}
\item $2.828$
\item $1.131$
\item $1.414$
\item $2.000$
\end{alternativas}
\gabarito{C}
\begin{sol}
\passo{steady state}{$k^*=\big(s/(n+\delta)\big)^{1/(1-\alpha)}=2^{3/2}=2.828$;
$y^*=(k^*)^{1/3}=2^{1/2}=1.414$. Option A is $k^*$; option B is $c^*=(1-s)y^*$.}
\refbib{Kurlat (2020), \S4.2, p.~57.}
\end{sol}

\medskip
\textbf{b)} (3 points) Moving the economy to the Golden Rule saving rate raises steady-state
consumption per worker by approximately:
\begin{alternativas}
\item $7.6\%$
\item $0\%$, since $s=0.20$ already maximizes $c^*$
\item $33.3\%$
\item $16.7\%$
\end{alternativas}
\gabarito{A}
\begin{sol}
\passo{Golden Rule}{$f'(k_{gr})=n+\delta$ gives $s_{gr}=\alpha=1/3$ and
$k_{gr}=(\alpha/(n+\delta))^{3/2}=6.086$, $c_{gr}=6.086^{1/3}-0.1\cdot6.086=1.217$.}
\passo{gain}{$c^*=0.8\cdot1.414=1.131$; $1.217/1.131-1=7.6\%$.}
\refbib{Kurlat (2020), \S4.3, p.~61.}
\end{sol}

\medskip
\textbf{c)} (2 points, T or F) Because this economy is below the Golden Rule, raising $s$ to
$s_{gr}$ is a Pareto improvement.
\gabarito{F}
\begin{sol}
\passo{transition}{Below the Golden Rule, reaching $k_{gr}$ requires \emph{lower} consumption
during the transition. Current generations lose; only the reverse case
($k^*>k_{gr}$, dynamic inefficiency) gives a free lunch.}
\refbib{Kurlat (2020), \S4.3.}
\end{sol}

\medskip
\textbf{d)} (2 points, T or F) If $g_Y=3\%$, $g_K=4.5\%$, $g_L=1.5\%$ and $\alpha=1/3$, then all
growth of output per worker comes from capital deepening.
\gabarito{F}
\begin{sol}
\passo{accounting}{$g_y=3-1.5=1.5\%$; capital deepening $\alpha g_k=\tfrac13(4.5-1.5)=1.0\%$;
residual $g_A=0.5\%$. A third of $g_y$ is TFP.}
\refbib{Kurlat (2020), \S5.4, eq.\ (5.4.2), p.~86.}
\end{sol}

% -----------------------------------------------------------------------------
\questao{(15 points) Discursive. Brazil is modeled as a Solow economy with labor-augmenting
technology: $Y=K^{\alpha}(AL)^{1-\alpha}$, $\alpha=1/3$, $s=0.24$, $n=0.01$, $g=0.02$ (growth of $A$),
$\delta=0.05$. Work \emph{per efficiency unit}: $\tilde k=K/(AL)$, $\tilde y=Y/(AL)$.}

\textbf{a)} (5 points) Derive the law of motion of $\tilde k$ (continuous time) and compute
$\tilde k^*$, $\tilde y^*$, $\tilde c^*$, $K/Y$ and the marginal product of capital in steady state.
Give the growth rates of $\tilde y$, $Y/L$ and $Y$ on the balanced growth path. Is Brazil above or
below the Golden Rule? Prove it with the marginal product, then confirm with $\tilde c$.
\resolespaco
\begin{sol}
\passo{law of motion}{$\dot{\tilde k}/\tilde k=\dot K/K-g-n$ and $\dot K=sY-\delta K$, so
$\dot{\tilde k}=s\tilde k^{\alpha}-(n+g+\delta)\tilde k$.}
\passo{steady state}{$\tilde k^*=(0.24/0.08)^{3/2}=3^{3/2}=5.196$, $\tilde y^*=3^{1/2}=1.732$,
$\tilde c^*=0.76\cdot1.732=1.316$, $K/Y=\tilde k^*/\tilde y^*=s/(n+g+\delta)=3$.}
\passo{growth}{$\tilde y$: 0; $Y/L$: $g=2\%$; $Y$: $n+g=3\%$.}
\passo{Golden Rule test}{$f'(\tilde k^*)=\alpha\tilde y^*/\tilde k^*=\tfrac13\cdot\tfrac13=0.111>n+g+\delta=0.08$.
Capital is scarce relative to the Golden Rule: \emph{below} it. Check: $s_{gr}=1/3$ gives
$\tilde k_{gr}=(0.333/0.08)^{3/2}=8.505$ and $\tilde c_{gr}=1.361>1.316$.}
\intuicao{The return on capital net of what it costs to maintain $\tilde k$ ($n+g+\delta$) is positive,
so more capital would raise steady-state consumption.}
\refbib{Kurlat (2020), \S4.3 and \S4.5, pp.~61--72.}
\rubrica{
\item Law of motion with $(n+g+\delta)$: 1 pt. Omitting $g$: max.\ 2 pts in the item.
\item Steady-state values and $K/Y$: 2 pts.
\item BGP growth rates: 1 pt. Saying $\tilde y$ grows at $g$: $-1$ pt.
\item Golden Rule comparison with a correct criterion: 1 pt.
}
\end{sol}

\textbf{b)} (5 points) Floods destroy 30\% of the capital stock at $t_0$ (illustrative). Using the
log-linear approximation $\dot{\ln\tilde y}\approx-\lambda(\ln\tilde y-\ln\tilde y^*)$, derive
$\lambda$, compute the half-life of the gap, the impact effect on $Y/L$ and its growth rate just
after the shock. Separately, suppose emigration permanently lowers $n$ to $0.005$: compute the
change in $\tilde y^*$ and the long-run growth of $Y/L$ and of $Y$.
\resolespaco
\begin{sol}
\passo{speed}{Linearize $\dot{\tilde k}/\tilde k=s\tilde k^{\alpha-1}-(n+g+\delta)$ in $\ln\tilde k$
at the steady state: derivative $(\alpha-1)s\tilde k^{*\alpha-1}=-(1-\alpha)(n+g+\delta)$. Since
$\ln\tilde y=\alpha\ln\tilde k$, the same speed applies:
$\lambda=(1-\alpha)(n+g+\delta)=\tfrac23\cdot0.08=0.0533$.}
\passo{half-life}{$\ln2/\lambda=13.0$ years.}
\passo{impact}{$L$ and $A$ do not jump, so $Y/L$ falls by $1-0.7^{1/3}=11.2\%$.}
\passo{growth after}{$\ln\tilde y-\ln\tilde y^*=\tfrac13\ln0.7=-0.119$; $\tilde y$ grows at
$0.0533\cdot0.119=0.63\%$, so $Y/L$ grows at $2+0.63=2.63\%$ and decelerates toward $2\%$.}
\passo{lower $n$}{$\tilde y^*=(0.24/0.075)^{1/2}=1.789$: $+3.3\%$ level effect. Long-run growth of $Y/L$
stays $g=2\%$; growth of $Y$ falls from $3\%$ to $2.5\%$.}
\intuicao{Destroyed capital raises its marginal product, so recovery is fast at first. Demography
changes the level of $\tilde y^*$ and the growth of the aggregate, never the growth of income per
person, which is pinned by $g$.}
\refbib{Kurlat (2020), \S4.2 and \S4.5; Romer (2012), \S1.5 (speed of convergence).}
\rubrica{
\item Derives $\lambda=(1-\alpha)(n+g+\delta)$: 1.5 pts. Quoting it without derivation: 0.5 pt.
\item Half-life and impact effect: 1.5 pts.
\item Growth just after the shock: 1 pt.
\item Effect of lower $n$, both growth rates: 1 pt. Saying long-run growth of $Y/L$ rises: max.\ 3 pts in the item.
}
\end{sol}

\textbf{c)} (5 points) Development accounting with $y=A\,k^{\alpha}h^{1-\alpha}$. Brazil/US ratios:
$y=0.25$, $k=0.20$, $h=0.60$. Compute the TFP ratio and the share of the log income gap due to
$k$, $h$ and TFP. Then suppose a quarter of Brazil's measured capital is security capital
(walls, armored cars) that produces nothing. Recompute the TFP ratio on productive capital and
explain what the residual measured in the first step was capturing.
\resolespaco
\begin{sol}
\passo{TFP}{$A_{BR}/A_{US}=0.25/(0.2^{1/3}\cdot0.6^{2/3})=0.25/(0.585\cdot0.711)=0.601$.}
\passo{shares}{$\ln0.25=-1.386$. Capital: $\tfrac13\ln0.2/\ln0.25=38.7\%$; human capital:
$\tfrac23\ln0.6/\ln0.25=24.6\%$; TFP: $\ln0.601/\ln0.25=36.7\%$. They add to $100\%$.}
\passo{security capital}{Productive $k$ ratio $=0.75\cdot0.2=0.15$:
$A_{BR}/A_{US}=0.25/(0.15^{1/3}\cdot0.711)=0.661$.}
\passo{interpretation}{Counting security capital as productive overstates Brazil's inputs, so
the measured residual (0.601) is lower than the TFP of productive capital (0.661). The residual
captures crime as if it were bad technology: it bundles technology with misallocation and
institutions.}
\intuicao{TFP is ``a measure of our ignorance'': whatever inputs do not explain lands in it,
including policies and resource allocation.}
\refbib{Kurlat (2020), \S5.3--5.5, pp.~80--95.}
\rubrica{
\item TFP ratio: 1.5 pts.
\item Log decomposition adding to one: 1.5 pts.
\item Recomputed TFP with productive capital: 1 pt.
\item Interpretation of the residual: 1 pt. Calling the residual ``pure technology'': 0 for this line.
}
\end{sol}

\end{document}
